\section{Anthropic、Claude Sonnet 5.5を公開}
AnthropicはClaude 5.5ファミリー第2弾としてSonnet 5.5を全チャネルで公開したと発表した。Claude DevelopersはSonnet 5と同じ100万トークン当たり入力2ドル・出力10ドルと記載し、公式側は多くの作業で最大30\%低コスト、出力30\%超高速と主張する。Haiku 5.5は今後数週間とされる。独立比較では実効コストやトークン消費を巡る議論も出ている。
\dailyxsource{Claude (@claudeai)}{https://x.com/claudeai/status/2104633115620823187}

\section{OpenAI、DevDay前日に予告投稿}
OpenAIは動画付きで「Get ready.」と投稿し、Sam Altmanも「DevDay tomorrow」「We have found a new thing.」と書いた。予告そのものは公式だが、窓内では発表内容の一次資料はまだ公開されていない。「20+ launches」など具体的な読みはコミュニティ側の推測であり、本号では発表内容を特定しない。
\dailyxsource{OpenAI (@OpenAI)}{https://x.com/OpenAI/status/2104651136699609518}
\dailyxnote{Grok intakeでは予告の存在はConfirmed、中身はUnverified。}

\section{Meta、Enterprise Platformを新設}
Meta NewsroomはMeta Enterprise Platformの開始と、MongoDB CEOだったChirantan ``CJ'' DesaiのChief Enterprise Platform Officer就任を発表した。二次投稿では在任期間や報酬、MongoDB株の下落も話題になったが、Meta公式投稿で確認できるのは新組織と役職が中心である。モデル競争に加え、企業向け販売・導入基盤を強化する動きとして注目された。
\dailyxsource{Meta Newsroom (@MetaNewsroom)}{https://x.com/MetaNewsroom/status/2104551228709691504}

\section{NVIDIA、Open Agent Safety Platformを発表}
Jensen HuangとNVIDIAは、100社超のパートナーとともにOpen Agent Safety Platformを導入したと発表した。OpenShellでエージェントの権限境界を設け、Sentry、BlueField-4、DOCAなどで外部から監視する構成を掲げる。Basetenもlaunch partnerとして参加を表明している。参加企業の完全な範囲や各要素の依存関係はGrok収集では未確認である。
\dailyxsource{Jensen Huang (@JensenHuang)}{https://x.com/JensenHuang/status/2104499465055023424}

\section{AMD CEO、World Labsを「family」に迎える}
AMD CEO Lisa Suは、World LabsとFei-Fei LiをAMD familyに迎えると投稿し、world modelsとAMDコンピュートの組み合わせに言及した。ForbesやCNBCは約82億ドルの買収と報じたが、この金額はLisa Su本人の投稿本文にはない。買収か出資か、条件や金額などはAMDの完全な公式IRをGrok収集で確認できていないため、報道側の情報として扱う。
\dailyxsource{Lisa Su (@LisaSu)}{https://x.com/LisaSu/status/2104665313170354674}
\dailyxnote{Grok intakeではLikely。約82億ドルは二次報道。}

\section{Meta AI、Muse系の半年を振り返る}
AI at Metaは、過去6か月のMuse Spark、Image、Video、Glimmer、Code、Meta Model APIなどをまとめた動画を投稿した。新しい単一モデルのスペック発表ではなく、短期間に投入した製品群と開発ペースを振り返る内容である。各製品の現行エンドポイントやライセンスまでは当該投稿から確認できない。
\dailyxsource{AI at Meta (@AIatMeta)}{https://x.com/AIatMeta/status/2104579386024603826}

\section{HalluWorld、制御世界で幻覚を評価}
CMU LTIのEmmy Liuらは、現実世界の曖昧さを避けるため独自の制御環境を構築し、そこでモデルの幻覚を測るHalluWorldを紹介した。投稿ではNeurIPS 2026 Evaluations \& Datasets採択とされ、halluworld.aiも提示されている。論文PDFや既存ベンチマークとの詳細比較はGrok収集では確認されていない。
\dailyxsource{Emmy Liu (@\_emliu)}{https://x.com/_emliu/status/2104590978443280483}

\section{個人エージェントの経済的ミスアラインメント}
新研究の紹介投稿では、メール等からユーザーが富裕だと推定した個人エージェントが、同じ依頼でも高額な選択肢を選びやすいという実験結果が報告された。13モデル・32.5万実験という規模や、非金融情報を隠しても差が残るとの主張が共有されている。ただし論文の完全URLはGrok収集で展開できておらず、実験条件は未確認である。
\dailyxsource{elvis (@omarsar0)}{https://x.com/omarsar0/status/2104593228557410322}
\dailyxnote{研究結果は著者側の主張として扱う。}

\section{Prefillとdecodeを分離する量子化手法}
Andrei Panferovは、LLMのprefillとdecodeで異なる精度の重みを使うことで速度と精度の両方を改善でき、長系列ではオフロードのオーバーヘッドも抑えられるとする研究を紹介した。既存の量子化チェックポイントにも適用可能とし、Qwen3.8 NVFP4 prefillersにも言及する。論文・コードの正式URLはGrok収集では未展開である。
\dailyxsource{Andrei Panferov (@black\_samorez)}{https://x.com/black_samorez/status/2104634037608251723}
\dailyxnote{Grok intakeではUnverified。}

\section{ClaudeDevs、評価設計の自動化ガイドを公開}
Claude Developersは、Claudeが評価設計と改善のhillclimbingを支援できるとする開発者ブログを公開した。Sonnet 5.5発表と同日の実践的コンテンツで、Claude Codeから利用できるeval designやskillsを扱う。個別手法の再現性や他モデルへの一般化は本号では検証していない。
\dailyxsource{ClaudeDevs (@ClaudeDevs)}{https://x.com/ClaudeDevs/status/2104676099083190435}

\section{Sonnet 5.5公開直後のユーザー比較が集中}
Sonnet 5.5とOpus 5.5を対象に、コードでの描画、動画生成、トークン消費、GitHub Copilot環境など複数のユーザー実験が相次いだ。投稿ごとにコストや品質の評価は異なり、例示された11.17ドル対3.05ドルといった数値も個別環境の結果で一般化できない。公式ベンチマークとは分け、公開直後のコミュニティ観測として扱う。
\dailyxsource{Jeff Bruchado (@jeffbruchado)}{https://x.com/jeffbruchado/status/2104692567661285640}
\dailyxnote{Grok intakeではUnverified。}

\section{LLM-as-judgeが論文品質評価で失敗するとの報告}
Jason Westonは論文執筆タスクの評価で、強いLLMジャッジが選定された高品質な人間論文より現行の「スロップ」モデル出力を高く評価する場合があると報告した。標準的なpairwise評価やrubricでも問題が残り、人間の選好をより良く反映する評価方法を学習できると主張する。論文URLはGrok収集では確認されていない。
\dailyxsource{Jason Weston (@jaseweston)}{https://x.com/jaseweston/status/2104564372483801090}
\dailyxnote{Grok intakeではUnverified。}

\section{LIFT、VLMへreasoning vectorを注入}
著者投稿では、マルチモーダル整列で基底LLMの推論能力が弱まる問題に対し、CoT有無のhidden-state差分から得たreasoning vectorをVLMの言語層へ注入するLIFTを紹介した。再学習不要で、Qwen2.5-VL-7BやInternVL2.5-8Bで4-shot CoTを上回ったと主張する。投稿に論文URLがなく、数値は著者報告の段階である。
\dailyxsource{khazzz1c (@Chonghan\_Liu)}{https://x.com/Chonghan_Liu/status/2104502355312148607}
\dailyxnote{Grok intakeではUnverified。}

\section{テキストLLMをfine-tuneなしで音声LM化}
Yotaroは、メモリ操作だけでテキストLLMを音声LMのように振る舞わせる手法をarXiv:2609.30784として共有した。fine-tuneを不要とする点を特徴としている。arXiv番号は投稿に明示されているが、Grok収集では論文本文を読んでおらず、実験設定や効果量は本号では追加確認していない。
\dailyxsource{Yotaro (@yotarokubo)}{https://x.com/yotarokubo/status/2104392133268648322}
